\documentclass[11pt]{article}
\usepackage[T1]{fontenc}
\usepackage{lmodern,amsmath,amssymb,amsthm,mathtools}
\usepackage[margin=1in]{geometry}
\usepackage{microtype,tikz-cd}
\usepackage[colorlinks=true,linkcolor=blue,citecolor=blue,urlcolor=blue]{hyperref}
\hypersetup{pdftitle={Maximality of the Dosen--Petric Boolean Category},pdfauthor={Florian Lengyel}}
\newtheorem{theorem}{Theorem}[section]
\newtheorem{lemma}[theorem]{Lemma}
\newtheorem{proposition}[theorem]{Proposition}
\newtheorem{corollary}[theorem]{Corollary}
\newcommand{\B}{\mathsf B}
\newcommand{\Occ}{\operatorname{Occ}}
\newcommand{\Rel}{\mathbf{Rel}_{\mathrm{fin}}}
\newcommand{\FormulaSet}{\operatorname{Form}}
\title{Maximality of the Do\v{s}en--Petri\'c Boolean Category}
\author{Florian Lengyel\thanks{Email: \href{mailto:flengyel@gradcenter.cuny.edu}{\texttt{flengyel@gradcenter.cuny.edu}}}}
\date{October 23, 2026}
\begin{document}
\maketitle

\begin{abstract}
The freely generated zero-identity Boolean category of Do\v{s}en and
Petri\'c is maximal: every additional equation between distinct parallel
arrows, closed under the proof operations and uniform formula
substitution, makes the quotient a preorder. This disproves the
nonmaximality conjecture in Section~14.2 of the September 2007 revision
of \emph{Proof-Theoretical Coherence}. Boolean Coherence reduces a new
equation to a contextual equation involving one literal. Uniform
substitution extends that equation to every formula. A symmetric arrow
constructed from the distributor then forces the two product
projections to coincide. Tagged occurrence sets give explicit
coordinates for the reduction and substitution.
\end{abstract}

\medskip
\noindent\textit{Mathematics Subject Classification (2020).}
Primary 03G30; Secondary 03F07, 18A32.

\smallskip
\noindent\textit{Keywords.} Boolean categories, maximality, categorial proof theory,
coherence, substitution congruences.

\section{Introduction}

Do\v{s}en and Petri\'c's Boolean proof category admits distinct arrows
with the same source and target. Their Boolean Coherence theorem
distinguishes these arrows by relations between occurrences of
propositional variables. In \cite[\S14.2, p.~322]{DP}, they leave open
whether the category is maximal and conjecture that it is not. We prove
maximality: adjoining any new arrow equation as a substitution-closed
schema identifies every pair of parallel arrows.

The proof has two parts. First, a relation distinguishing two arrows
contains an edge that can be isolated by composition with a selector
and a coselector. This yields the contextual equation used in
Do\v{s}en and Petri\'c's relative-maximality argument
\cite[\S9.7, pp.~215--216]{DP}. Second, the distributor and zero
identities turn that contextual equation into equality of the product
projections. This second step proves the required collapse.

Do\v{s}en and Petri\'c prove maximality for cartesian categories
\cite[\S5]{DPCartesian}. Their proof uses coherence, uniform
substitution, and precomposition and postcomposition to reduce a new
cartesian equation to equality of the product projections. The
corresponding reduction for Boolean arrow equations is proved in
Sections~\ref{sec:isolation} and~\ref{sec:collapse}.

Earlier results establish maximality for lattice categories and relative
maximality for dicartesian categories
\cite[\S\S8--10]{DPBicartesian}; these results concern products and
coproducts without distribution. Related occurrence-selection
constructions appear in \cite[Lemma~6]{DPIsomorphic}. A differing-edge
argument also proves maximality for presentations of split preorders,
split equivalence relations, and binary relations \cite[\S16]{DPSplit}.

Lamarche and Stra{\ss}burger give a proof-net construction of free
categories for their Boolean-category axioms
\cite[Theorem~3.5.4]{LSFree}.
Stra{\ss}burger studies axioms for $*$-autonomous categories with
monoid and comonoid structures \cite{Strassburger}.
F\"uhrmann and Pym interpret cut reduction by an order on arrows and
prove soundness and completeness for their net theories
\cite[Theorems~6.4--6.5]{FP}.
After precomposition with the inverse right unit isomorphism, their
negation triangle \cite[(11), p.~28]{FP} equates an endomorphism with
the identity. The corresponding structural composite defines the
zero identity in \eqref{eq:zero-identity}.

We use the occurrence coordinates of \cite[\S9]{Notation}: occurrences
in a compound formula form tagged disjoint unions, and substitution
replaces each occurrence by its own block of occurrences.
Section~\ref{sec:occurrences} describes these blocks, and
Appendix~\ref{app:substitution} gives the commutative diagrams relating
them to Do\v{s}en and Petri\'c's ordinal substitution operation.
References to \cite{DP} use the printed pagination of its September
2007 revised version.

\section{The Boolean proof category}\label{sec:category}

Fix a nonempty set $P$ of propositional variables. Write $\FormulaSet(P)$
for the formulas built from $P$, the constants $\top,\bot$, the binary
connectives $\wedge,\vee$, and negation $\neg$. Let $\B(P)$ be the
freely generated zero-identity Boolean category defined in
\cite[\S\S14.2--14.3]{DP}. Its objects are these formulas, and its
arrows are the arrow terms of that presentation modulo its equations.
We write $\B$ when $P$ is fixed. We recall the structure and equations
used in the proof.

The connectives $\wedge$ and $\vee$ give binary products and coproducts;
$\top$ is terminal and $\bot$ is initial. For $A_1=A$ and $A_2=B$,
denote the projections and injections by
\[
 \pi_i^{A,B}:A\wedge B\longrightarrow A_i,
 \qquad
 \iota_i^{A,B}:A_i\longrightarrow A\vee B
 \qquad(i=1,2).
\]
Pairing and copairing are denoted
by $\langle f,g\rangle$ and $[f,g]$, respectively. Write
$!_A:A\to\top$ and $\kappa_A:\bot\to A$ for the terminal and initial
arrows, and $1_A$ for the identity of $A$. Composition $g\circ f$
applies $f$ first. The symbols $f\wedge g$ and $f\vee g$ denote the
arrows induced by the product and coproduct bifunctors.

The distributor is an arrow
\[
 d_{A,B,C}:A\wedge(B\vee C)\longrightarrow(A\wedge B)\vee C.
\]
Its first-injection equation is
\begin{equation}
 d_{A,B,C}\circ(1_A\wedge\iota_1^{B,C})
   =\iota_1^{A\wedge B,C}.
 \label{eq:first-injection}
\end{equation}
This is the coproduct-injection equation of
\cite[\S11.1, p.~234]{DP}.

Negation is represented by a functor
$\neg:\B^{\mathrm{op}}\to\B$, with specified natural
double-negation and De Morgan isomorphisms. In particular, write
$n_A:\neg\neg A\to A$ for the double-negation isomorphism.
The presentation also has arrows
$\eta_A:\top\to\neg A\vee A$ and
$\epsilon_A:A\wedge\neg A\to\bot$.
The \emph{zero identity} is the endomorphism $0_A:A\to A$ defined by
\begin{equation}
 0_A=[\kappa_A,1_A]\circ(\epsilon_A\vee1_A)\circ d_{A,\neg A,A}
     \circ(1_A\wedge\eta_A)\circ\langle1_A,!_A\rangle.
 \label{eq:zero-identity}
\end{equation}
The defining zero-identity equation asserts that
\begin{equation}
 f\circ0_A=0_C\circ g
 \qquad(f,g:A\to C).
 \label{eq:zero-law}
\end{equation}
These are the equations of \cite[\S14.2, p.~317]{DP}, with the
product and coproduct unit isomorphisms written explicitly.

\noindent\textit{Remark.}
Appendix~\ref{app:structural-diagrams} gives diagrammatic forms of
equations~\eqref{eq:first-injection}--\eqref{eq:zero-law}, with the
source and target of every arrow displayed.

Whenever $\B(A,C)$ is nonempty, define $0_{A,C}=f\circ0_A$ for any
$f:A\to C$. Equation~\eqref{eq:zero-law} makes this independent of
$f$: both $f\circ0_A$ and $g\circ0_A$ equal $0_C\circ g$.
Taking $f=g$ gives $f\circ0_A=0_C\circ f$. Consequently zero arrows
absorb composition whenever the composites are defined:
$h\circ0_{A,C}=0_{A,D}$ for $h:C\to D$, and
$0_{A,C}\circ k=0_{X,C}$ for $k:X\to A$.

Two arrows are \emph{parallel} when they have the same source and
target. A \emph{proof congruence} on $\B$ is an equivalence relation
on each hom-set preserved by composition and every arrow operation
of the presentation, including pairing, copairing, and negation.
A \emph{substitution congruence} is a proof congruence also preserved
by every uniform substitution $\sigma:P\to\FormulaSet(P)$.
Here substitution acts on the entire arrow term, including its source
and target; this is rule~$(su)$ of \cite[\S2.7, p.~51]{DP}.

The \emph{total congruence} identifies every parallel pair. A category
is \emph{thin}, or a \emph{preorder}, when each hom-set has at most one
element. We call $\B$ \emph{maximal} when its only substitution
congruences are equality and the total congruence. The main result is
the following theorem.

\begin{theorem}\label{thm:main}
Let $P$ be a nonempty set of propositional variables. The freely
generated zero-identity Boolean category $\B(P)$ of Do\v{s}en and
Petri\'c is maximal. Equivalently, every substitution congruence on
$\B(P)$ that identifies two distinct parallel arrows identifies all
parallel arrows.
\end{theorem}

The proof is completed in Section~\ref{sec:collapse}.

\section{Occurrences and substitution}\label{sec:occurrences}

For a formula $A$, write $\Occ(A)$ for the finite set of positions
occupied by propositional variables in $A$.
Appendix~\ref{app:profiles} describes substitution on these occurrence
sets by commutative diagrams. Repeated occurrences of the same variable
occupy distinct positions. We use tagged coordinates for these positions.

For finite sets $X,Y$, let
$X\sqcup Y=(\{1\}\times X)\cup(\{2\}\times Y)$ be their tagged
disjoint union. The occurrence sets are defined recursively by
\[
 \begin{aligned}
 \Occ(p)&=\{*\},&
 \Occ(\top)&=\Occ(\bot)=\varnothing,\\
 \Occ(\neg A)&=\Occ(A),&
 \Occ(A\mathbin\star C)&=\Occ(A)\sqcup\Occ(C)
       \quad(\star\in\{\wedge,\vee\}).
 \end{aligned}
\]
The label map $\lambda_A:\Occ(A)\to P$ assigns to each occurrence
its underlying variable. Distinct positions may have the same label.
The left-to-right
order gives a bijection from $\Occ(A)$ to
$\{1,\ldots,|\Occ(A)|\}$, interpreted as the empty set when the
cardinality is zero. This bijection relates the tagged coordinates to
the finite ordinals used in \cite{DP}; Appendix~\ref{app:ordinals}
expresses the corresponding change of substitution coordinates by a
commutative square and a worked example.

Let $\Rel$ be the category of finite sets and binary relations.
A relation $R:X\to Y$ is a subset of $X\times Y$; identities are
diagonal relations, and composition is relational composition.
Transporting the occurrence interpretation of \cite{DP} to the
tagged coordinates gives a functor
$G:\B\to\Rel$ with $GA=\Occ(A)$.
For example,
\[
 G\pi_1^{A,C}=\{((1,x),x):x\in\Occ(A)\},\qquad
 G\iota_1^{A,C}=\{(x,(1,x)):x\in\Occ(A)\}.
\]
The second projection and injection use tag~$2$. Pairing takes the
union of the two relations after tagging their targets; copairing
takes the union after tagging their sources. The relation $Gd_{A,B,C}$
is the graph of the rebracketing bijection
\[
 \Occ(A)\sqcup\bigl(\Occ(B)\sqcup\Occ(C)\bigr)
 \longrightarrow
 \bigl(\Occ(A)\sqcup\Occ(B)\bigr)\sqcup\Occ(C),
\]
which preserves each occurrence of $A,B,C$
\cite[\S2.9, p.~60]{DP}. Zero arrows have empty relations.

\begin{theorem}[Boolean Coherence, Do\v{s}en--Petri\'c]
\label{thm:coherence}
The functor $G:\B\to\Rel$ is faithful: parallel arrows $f,g$ of
$\B$ are equal if and only if $Gf=Gg$.
\end{theorem}

This is \cite[\S14.2, p.~321]{DP}. In particular, for $p\in P$, the
two projections $p\wedge p\to p$ have different occurrence
relations, so $\B$ is nonthin. Below, coherence is applied to
equations in $\B$; these equations then pass to its quotients.

A \emph{signed literal} is $p$ or $\neg p$ for $p\in P$.
A formula is in \emph{negation normal form} when negation occurs only
immediately before a propositional variable. The double-negation and
De Morgan isomorphisms, including the isomorphisms for the truth
constants, give an isomorphism $a_A:A\to A^\circ$ to a formula in
negation normal form. Its occurrence relation is the graph of the
correspondence of positions \cite[\S14.1, pp.~313--315]{DP}.

For arrows between formulas in negation normal form, every edge of
the occurrence relation joins occurrences of the same signed literal.
Indeed, the equivalent presentation of \cite[\S14.2, pp.~318--321]{DP}
has signed literals as generators. Its lattice generators preserve
these labels; the arrows $\eta_p$ and $\epsilon_p$ have empty
relations. Composition, pairing, and copairing preserve label
agreement. The equivalence with $\B$ preserves the occurrence
interpretation.

Substitution replaces each occurrence of a variable by a separate
block of occurrences of its substituted formula.
Appendix~\ref{app:profiles} describes these blocks and their label maps
by commutative diagrams, following \cite[\S9]{Notation}.
Let $Q$ be a set of propositional variables and
$\sigma:P\to\FormulaSet(Q)$ a substitution. If $A[\sigma]$ denotes the
result of uniform substitution, there is a canonical bijection
\begin{equation}
 \Occ(A[\sigma])\cong D_{A,\sigma},\qquad
 D_{A,\sigma}=\coprod_{x\in\Occ(A)}
                  \Occ\bigl(\sigma(\lambda_A(x))\bigr).
 \label{eq:occ-substitution}
\end{equation}
The set on the right consists of pairs $(x,t)$, where $x$ is an old
occurrence and $t$ is an occurrence within the formula inserted there.
Induction on $A$ proves the bijection: the variable case has one
block, the constant cases have none, binary connectives form tagged
unions, and negation preserves the occurrence set. In particular,
two occurrences of the same variable give two separately tagged
copies of its substituted formula.

The organization of these blocks is expressed by a square of sets and
functions. Put $E_\sigma=\coprod_{p\in P}\Occ(\sigma(p))$.
Define $\operatorname{pr}_D:D_{A,\sigma}\to\Occ(A)$ and
$\operatorname{pr}_E:E_\sigma\to P$ by
$\operatorname{pr}_D(x,t)=x$ and $\operatorname{pr}_E(p,t)=p$.
The function
$\widehat\lambda_{A,\sigma}:D_{A,\sigma}\to E_\sigma$ sends
$(x,t)$ to $(\lambda_A(x),t)$. Thus
\begin{equation}
\begin{tikzcd}[column sep=large,row sep=large]
 D_{A,\sigma}\arrow[r,"\widehat\lambda_{A,\sigma}"]
   \arrow[d,"\operatorname{pr}_D"']
   &E_\sigma\arrow[d,"\operatorname{pr}_E"]\\
 \Occ(A)\arrow[r,"\lambda_A"']&P.
\end{tikzcd}
\label{eq:block-square}
\end{equation}
Both routes send $(x,t)$ to the variable $\lambda_A(x)$ whose
substituted formula contains $t$. The upper arrow forgets which
occurrence of that variable supplied the copy; the source $D_{A,\sigma}$
retains it.

Substitution on arrow terms induces a functor
\[
 F_\sigma:\B(P)\longrightarrow\B(Q),\qquad
 A\longmapsto A[\sigma],\quad f\longmapsto f[\sigma].
\]
Rule~$(su)$ makes this functor well defined on arrow terms modulo the
defining equations. It preserves identities, composition, and the
specified arrow operations. For each arrow, the blocks in
\eqref{eq:occ-substitution} describe the occurrence sets of its
substituted source and target.

\section{Isolating an occurrence}\label{sec:isolation}

Put $j_A=\iota_1^{A,\top}:A\to A\vee\top$. The contextual equation
needed in the proof is
\begin{equation}
 j_A\circ\pi_1^{A,\bot}
   =0_{A\wedge\bot,A\vee\top}.
 \tag{$K_A$}\label{eq:context}
\end{equation}
The right-hand side is also
$\iota_2^{A,\top}\circ!_\bot\circ\pi_2^{A,\bot}$: its occurrence
relation is empty, so Boolean Coherence identifies it with the zero
arrow. Thus \eqref{eq:context} is the contextual equation of
\cite[\S9.7, pp.~215--216]{DP}.

\begin{lemma}[Occurrence isolation]\label{lem:isolation}
For distinct parallel arrows $f,g:A\to C$ in $\B$, there are a
signed literal $\ell$ and arrows
\[
 i:\ell\wedge\bot\longrightarrow A,
 \qquad r:C\longrightarrow\ell\vee\top
\]
such that, after possibly interchanging $f$ and $g$,
\[
 r\circ f\circ i=j_\ell\circ\pi_1^{\ell,\bot},
 \qquad
 r\circ g\circ i=0_{\ell\wedge\bot,\ell\vee\top}.
\]
\end{lemma}

\begin{proof}
First replace $f,g$ by $a_C\circ f\circ a_A^{-1}$ and
$a_C\circ g\circ a_A^{-1}$. These remain distinct, so it suffices to
construct the arrows for the negation normal forms $A^\circ,C^\circ$.
If their selector and coselector are $i^\circ,r^\circ$, the required
arrows for the original formulas are
$a_A^{-1}\circ i^\circ$ and $r^\circ\circ a_C$.
We may therefore assume that $A,C$ are in negation normal form.

Faithfulness of $G$ gives an edge belonging to exactly one of $Gf,Gg$.
Interchange $f,g$ if necessary, and choose
$(x,y)\in Gf\setminus Gg$. The occurrences $x$ and $y$ have the same
signed literal $\ell$. Set $D=\ell\wedge\bot$ and $E=\ell\vee\top$.
For every formula $V$, the arrows
\[
 b_V=\kappa_V\circ\pi_2^{\ell,\bot}:D\to V,
 \qquad
 t_V=\iota_2^{\ell,\top}\circ!_V:V\to E
\]
have empty occurrence relations.

Construct the selector $i:D\to A$ recursively along the tagged path
to $x$. At the selected literal, use $\pi_1^{\ell,\bot}$.
At a conjunction, pair the recursively constructed arrow with $b_V$
for the unselected factor, placing it in the order specified by the
tag. At a disjunction, compose the recursively constructed arrow
with the injection of the selected summand. The displayed rules for
$G$ show inductively that $Gi$ consists of the single edge from the
sole occurrence of $D$ to $x$.

Construct $r:C\to E$ dually along the path to $y$.
At the selected literal, use $\iota_1^{\ell,\top}$.
At a conjunction, first project to the selected factor and then apply
the recursively constructed arrow. At a disjunction, copair that
arrow with $t_V$ on the unselected summand. Thus $Gr$ is the single
edge from $y$ to the sole occurrence of $E$.

Relational composition now makes $G(r\circ f\circ i)$ the one-edge
relation from $\Occ(D)$ to $\Occ(E)$, whereas
$G(r\circ g\circ i)$ is empty. Boolean Coherence gives the asserted
equations.
\end{proof}

\begin{proposition}\label{prop:all-contexts}
If a substitution congruence $\theta$ on $\B$ identifies distinct
parallel arrows, then \eqref{eq:context} holds modulo $\theta$ for
every formula $A$.
\end{proposition}

\begin{proof}
Lemma~\ref{lem:isolation} and compatibility with composition give
\eqref{eq:context} for a signed literal $\ell$.
If $\ell=p$, let $\sigma$ substitute $A$ for $p$ and fix the other variables.
This yields \eqref{eq:context} for $A$: under
\eqref{eq:occ-substitution}, the literal occurrence in each endpoint
becomes the block $\Occ(A)$, and the arrow operations are preserved
by $F_\sigma$. Appendix~\ref{app:shared} describes this
one-occurrence case of the substitution diagrams.

If $\ell=\neg p$, substitute $\neg A$ for $p$, obtaining the equation
for $\neg\neg A$. Transport it along $n_A:\neg\neg A\to A$.
Explicitly, let
$k_A=j_A\circ\pi_1^{A,\bot}:A\wedge\bot\to A\vee\top$.
The product and coproduct equations give the commutative square
\[
\begin{tikzcd}[column sep=large,row sep=large]
 A\wedge\bot\arrow[r,"k_A"]\arrow[d,"n_A^{-1}\wedge1_\bot"']
   &A\vee\top\\
 \neg\neg A\wedge\bot\arrow[r,"k_{\neg\neg A}"']
   &\neg\neg A\vee\top\arrow[u,"n_A\vee1_\top"'].
\end{tikzcd}
\]
Precomposition and postcomposition of the zero arrow give a zero
arrow by absorption. The equation for $\neg\neg A$ therefore gives
\eqref{eq:context} for $A$.
\end{proof}

\section{Contextual collapse}\label{sec:collapse}

The distributor produces an arrow that treats two copies of $A$
symmetrically. This symmetry converts the contextual equation into
an equation between the projections.

\begin{lemma}\label{lem:symmetric}
For formulas $A,B$, define
\[
 \begin{aligned}
 u_A^B&=d_{A,B,A}\circ(1_A\wedge\iota_2^{B,A})
       :A\wedge A\longrightarrow(A\wedge B)\vee A,\\
 h_A^B&=[\pi_1^{A,B},1_A]\circ u_A^B:A\wedge A\longrightarrow A.
 \end{aligned}
\]
The arrow $h_A^B$ is independent of $B$. Writing it as $h_A$ and
putting $\tau_A=\langle\pi_2^{A,A},\pi_1^{A,A}\rangle$, one has
$h_A\circ\tau_A=h_A$.
\end{lemma}

\begin{proof}
The relation of $u_A^B$ sends $(1,x)$ to $(1,(1,x))$ and $(2,x)$
to $(2,x)$ in $\bigl(\Occ(A)\sqcup\Occ(B)\bigr)\sqcup\Occ(A)$.
Copairing with $\pi_1^{A,B}$ and $1_A$ then gives
\[
 Gh_A^B=\{((1,x),x),((2,x),x):x\in\Occ(A)\}.
\]
This relation is independent of $B$ and invariant under exchange of
the two source tags. Boolean Coherence proves both assertions.
\end{proof}

\begin{lemma}[Contextual collapse]\label{lem:collapse}
Let $\theta$ be a proof congruence on $\B$ such that
\eqref{eq:context} holds modulo $\theta$ for every formula $A$.
Then $\B/\theta$ is a preorder.
\end{lemma}

\begin{proof}
Write $\equiv$ for equality modulo $\theta$.
The quotient retains products and coproducts: their pairing,
copairing, and defining equations pass to equivalence classes.
Using $B=\bot$ in Lemma~\ref{lem:symmetric}, copairing and
\eqref{eq:context} give
\begin{equation}
 \begin{aligned}
 j_A\circ h_A
   &=[j_A\circ\pi_1^{A,\bot},j_A]\circ u_A^\bot\\
   &\equiv[0_{A\wedge\bot,A\vee\top},j_A]\circ u_A^\bot\\
   &=j_A\circ\pi_2^{A,A}.
 \end{aligned}
 \label{eq:second-block}
\end{equation}
For the last equality, both occurrence relations discard the first
source tag and send $(2,x)$ to $(1,x)$ in $\Occ(A)\sqcup\varnothing$.
Thus it is an equation of $\B$ by Boolean Coherence.
Precomposing \eqref{eq:second-block} with $\tau_A$ and using
$h_A\circ\tau_A=h_A$ yields
\[
 j_A\circ\pi_1^{A,A}\equiv j_A\circ\pi_2^{A,A}.
\]
Now precompose with $\langle1_A,0_A\rangle:A\to A\wedge A$.
It follows that
\begin{equation}
 j_A\equiv j_A\circ0_A
   =0_{A,A\vee\top}
   =\iota_2^{A,\top}\circ!_A.
 \label{eq:terminal-injection}
\end{equation}
The last equality holds by coherence, since the composite through
$\top$ has empty occurrence relation.

The arrows $s=\iota_2^{A,\top}:\top\to A\vee\top$ and
$t=!_{A\vee\top}:A\vee\top\to\top$ are inverse in the quotient.
Indeed, $t\circ s=1_\top$ by terminality. Equation
\eqref{eq:terminal-injection} gives
$(s\circ t)\circ\iota_1^{A,\top}\equiv\iota_1^{A,\top}$, and
terminality gives
$(s\circ t)\circ\iota_2^{A,\top}=\iota_2^{A,\top}$.
The coproduct property therefore gives $s\circ t\equiv1_{A\vee\top}$.
Hence $A\vee\top$, and by coproduct symmetry also $\top\vee A$,
is terminal in the quotient. In particular,
\begin{equation}
 \iota_2^{\top,A}\equiv\iota_1^{\top,A}\circ!_A.
 \label{eq:other-injection}
\end{equation}

Evaluate the same arrow $h_A$ using $B=\top$.
Substituting \eqref{eq:other-injection} into its defining expression
and applying \eqref{eq:first-injection} gives
\[
 \begin{aligned}
 h_A
   &\equiv[\pi_1^{A,\top},1_A]\circ d_{A,\top,A}
       \circ(1_A\wedge\iota_1^{\top,A})
       \circ(1_A\wedge!_A)\\
   &=[\pi_1^{A,\top},1_A]\circ\iota_1^{A\wedge\top,A}
       \circ(1_A\wedge!_A)\\
   &=\pi_1^{A,A}.
 \end{aligned}
\]
Since $h_A$ is invariant under $\tau_A$, this implies
$\pi_1^{A,A}\equiv\pi_2^{A,A}$.
For any parallel arrows $f,g:X\to A$, pairing now gives
\[
 f=\pi_1^{A,A}\circ\langle f,g\rangle
   \equiv\pi_2^{A,A}\circ\langle f,g\rangle=g.
\]
Thus the quotient is a preorder.
\end{proof}

\begin{proof}[Proof of Theorem~\ref{thm:main}]
Suppose that a substitution congruence $\theta$ on $\B(P)$ identifies
two distinct parallel arrows. Proposition~\ref{prop:all-contexts}
gives \eqref{eq:context} modulo $\theta$ for every formula $A$.
Lemma~\ref{lem:collapse} then makes $\theta$ total.
\end{proof}

The same argument gives the formulation of maximality in terms of
models used in \cite[\S9.3, p.~194]{DP}. An \emph{interpretation} in
a zero-identity Boolean category $\mathcal C$ assigns an object of
$\mathcal C$ to each propositional variable and interprets the
specified operations; it induces a structure-preserving functor
$\B(P)\to\mathcal C$.

\begin{corollary}\label{cor:models}
Let $P$ be a nonempty set of propositional variables, and let
$f,g:A\to C$ be distinct arrows of $\B(P)$.
If a zero-identity Boolean category $\mathcal C$ satisfies $f=g$
under every interpretation of the propositional variables, then
$\mathcal C$ is a preorder.
\end{corollary}

\begin{proof}
Declare two parallel arrows of $\B(P)$ equivalent when all their
interpretations in $\mathcal C$ agree. This defines a substitution
congruence: composing an interpretation with $F_\sigma$ is another
interpretation. The congruence contains $f=g$, so
Theorem~\ref{thm:main} makes it total. In particular,
$\pi_1^{p,p}=\pi_2^{p,p}$ holds under every interpretation, for any
$p\in P$. Interpret $p$ as an arbitrary object $X$ of $\mathcal C$.
Its two projections $X\wedge X\to X$ are equal. Pairing any arrows
$a,b:Y\to X$ with these projections yields $a=b$.
\end{proof}

\appendix
\section{Substitution in occurrence coordinates}\label{app:substitution}

The generalized composition of definable connectives in
\cite[\S2.5, pp.~45--46]{DP} records how repeated arguments are
substituted into repeated positions. We give its diagrammatic description
from \cite[\S\S9.2--9.3]{Notation}. The diagrams separate the choice of
an outer occurrence, the choice of an occurrence within an inserted
formula, and the name assigned to that occurrence.

\subsection{Profiles and their fibres}\label{app:profiles}

For an integer $n\geq0$, put $\underline n=\{1,\ldots,n\}$, with
$\underline0=\varnothing$. A \emph{formula pattern} on a finite set
of placeholders is a formula in those placeholders. Its
\emph{occurrence profile} records the placeholder at each position,
reading positions from left to right.

Let $U=\{p_1,\ldots,p_k\}$ consist of distinct placeholders and let
$M\in\FormulaSet(U)$ have $m$ occurrences. Its profile is
\[
 o_M:\underline m\longrightarrow\underline k,
 \qquad o_M(r)=i\quad\text{when position $r$ contains $p_i$}.
\]
For each $i\in\underline k$, choose a set
$Q_i=\{q_{i,1},\ldots,q_{i,\ell_i}\}$ of distinct local placeholders and a
pattern $N_i\in\FormulaSet(Q_i)$ with $n_i$ occurrences. The sets
$Q_i$ are pairwise disjoint. Write
$o_i:\underline{n_i}\to\underline{\ell_i}$ for the profile of $N_i$.
All these numbers may be zero whenever the displayed maps exist.

The occurrences after substitution and the local placeholder indices
form, respectively, the sets
\[
 \begin{aligned}
 D&=\{(r,t):r\in\underline m,\ t\in\underline{n_{o_M(r)}}\},\\
 C&=\{(i,s):i\in\underline k,\ s\in\underline{\ell_i}\}.
 \end{aligned}
\]
Thus $(r,t)$ specifies an occurrence in the copy of $N_{o_M(r)}$
inserted at position $r$ of $M$. The pair $(i,s)$ specifies a local
placeholder of $N_i$. Define the injections
\[
 a_r:\underline{n_{o_M(r)}}\longrightarrow D,\quad a_r(t)=(r,t),
 \qquad
 b_i:\underline{\ell_i}\longrightarrow C,\quad b_i(s)=(i,s).
\]
There is a unique function $\gamma:D\to C$ for which, for every
$r\in\underline m$, the following square commutes:
\begin{equation}\label{eq:app-fibre}
\begin{tikzcd}[column sep=large,row sep=large]
 \underline{n_{o_M(r)}}\arrow[r,"a_r"]\arrow[d,"o_{o_M(r)}"']
   &D\arrow[d,"\gamma"]\\
 \underline{\ell_{o_M(r)}}\arrow[r,"b_{o_M(r)}"']&C.
\end{tikzcd}
\end{equation}
The square says that an occurrence inside one inserted copy is labelled
by that pattern's profile and then placed in the corresponding target
block. Since $D$ is the disjoint union of its source blocks, these
squares determine $\gamma$ on every element. Explicitly,
\[
 \gamma(r,t)=\bigl(o_M(r),o_{o_M(r)}(t)\bigr).
\]

Define the block projections $p_D:D\to\underline m$ and
$p_C:C\to\underline k$ by $p_D(r,t)=r$ and $p_C(i,s)=i$.
The same construction fits over the outer profile:
\begin{equation}\label{eq:app-profile}
\begin{tikzcd}[column sep=large,row sep=large]
 D\arrow[r,"\gamma"]\arrow[d,"p_D"']&C\arrow[d,"p_C"]\\
 \underline m\arrow[r,"o_M"']&\underline k.
\end{tikzcd}
\end{equation}
Both routes send $(r,t)$ to $o_M(r)$. This square records which
inserted pattern is used; square~\eqref{eq:app-fibre} specifies the
map within each such block. Repeated values of $o_M$ give separate
source blocks, one for each occurrence of the corresponding outer
placeholder.

To identify this construction with actual pattern substitution,
replace each $q_{i,s}$ in $N_i$ by $(i,s)$ and denote the resulting
pattern over $C$ by $\widetilde N_i$. Let
$\Phi_0:U\to\FormulaSet(C)$ be given by $\Phi_0(p_i)=\widetilde N_i$.
Every occurrence of $M[\Phi_0]$ has unique coordinates $(r,t)\in D$.
This gives a bijection $\chi:\Occ(M[\Phi_0])\to D$, and its relation
to the label map is the commutative square
\begin{equation}\label{eq:app-occurrence-labels}
\begin{tikzcd}[column sep=large,row sep=large]
 \Occ(M[\Phi_0])\arrow[r,"\chi","\cong"']
    \arrow[d,"\lambda_{M[\Phi_0]}"']&D\arrow[d,"\gamma"]\\
 C\arrow[r,"1_C"']&C.
\end{tikzcd}
\end{equation}
Indeed, the occurrence with coordinates $(r,t)$ contains precisely
the placeholder $\bigl(o_M(r),o_{o_M(r)}(t)\bigr)$.

\subsection{The finite-ordinal presentation}\label{app:ordinals}

The finite ordinals used in \cite{DP} number the elements of each
block consecutively. Put
\[
 s_0=t_0=0,\qquad
 s_r=\sum_{q=1}^{r}n_{o_M(q)}\quad(1\leq r\leq m),\qquad
 t_i=\sum_{j=1}^{i}\ell_j\quad(1\leq i\leq k).
\]
The block encodings are the bijections
\[
 \begin{aligned}
 \theta_D:D&\longrightarrow\underline{s_m},
       &\theta_D(r,t)&=s_{r-1}+t,\\
 \theta_C:C&\longrightarrow\underline{t_k},
       &\theta_C(i,s)&=t_{i-1}+s.
 \end{aligned}
\]
Define $\delta:\underline{s_m}\to\underline{t_k}$ as the unique
function making the following square commute:
\begin{equation}\label{eq:app-ordinal}
\begin{tikzcd}[column sep=large,row sep=large]
 D\arrow[r,"\gamma"]\arrow[d,"\theta_D","\cong"']
    &C\arrow[d,"\theta_C","\cong"']\\
 \underline{s_m}\arrow[r,"\delta"']&\underline{t_k}.
\end{tikzcd}
\end{equation}
The vertical maps change coordinates. The top route keeps the block
and local position visible; the bottom route uses a single position
number. In coordinates, commutativity says
\begin{equation}\label{eq:app-coordinate}
 \delta(s_{r-1}+t)=t_{o_M(r)-1}+o_{o_M(r)}(t)
 \qquad\bigl((r,t)\in D\bigr).
\end{equation}
Every source position belongs to exactly one nonempty block, so this
formula determines $\delta$ even when some blocks are empty.

Do\v{s}en and Petri\'c write $\mu=o_M$ and $\nu_i=o_i$, and denote
$\delta$ by $\mu(\nu_1,\ldots,\nu_k)$. Their block locator selects the
unique $r$ for which $s_{r-1}<x\leq s_r$; substituting
$t=x-s_{r-1}$ in \eqref{eq:app-coordinate} gives exactly their
formula \cite[\S2.5, p.~46]{DP}. The diagram therefore describes the
same generalized composition. A pattern can equivalently be
specified by erasing its placeholder names to one common symbol and
attaching its occurrence profile. In that notation, the generalized
composite is represented by the skeleton of the substituted pattern
together with the profile $\delta$.

For example, take
\[
 M(p,q)=q\wedge(p\vee q),\qquad
 N_1(a,b)=a\wedge b,\qquad
 N_2(c,d)=d\wedge(c\vee d).
\]
The profiles, listed by successive values, are
$o_M=(2,1,2)$, $o_1=(1,2)$, and $o_2=(2,1,2)$.
The source blocks have sizes $3,2,3$, while the target blocks list
$(a,b)$ and $(c,d)$. Thus
\[
\begin{array}{c|ccc|cc|ccc}
 \text{source position}&1&2&3&4&5&6&7&8\\\hline
 \text{placeholder}&d&c&d&a&b&d&c&d\\
 \delta&4&3&4&1&2&4&3&4
\end{array}
\]
For instance, source position $2$ has coordinates $(1,2)$.
Square~\eqref{eq:app-fibre} sends it to $(2,1)$, and
$\theta_C(2,1)=3$. This is the entry $\delta(2)=3$.

\subsection{Shared variables and the maximality argument}\label{app:shared}

The local placeholders may be assigned names in a common variable
set $Q$. Let $\rho:C\to Q$ be that assignment; different elements
of $C$ may have the same image. Replace each placeholder $c\in C$
by the variable $\rho(c)$, obtaining the renaming map
$(-)[\rho]:\FormulaSet(C)\to\FormulaSet(Q)$. Define
$\sigma:U\to\FormulaSet(Q)$ by
$\sigma(p_i)=\widetilde N_i[\rho]$.
Then the following square commutes:
\begin{equation}\label{eq:app-renaming}
\begin{tikzcd}[column sep=large,row sep=large]
 \FormulaSet(U)\arrow[r,"{(-)[\Phi_0]}"]
    \arrow[d,"{(-)[\sigma]}"']
   &\FormulaSet(C)\arrow[d,"{(-)[\rho]}"]\\
 \FormulaSet(Q)\arrow[r,"1"']&\FormulaSet(Q).
\end{tikzcd}
\end{equation}
At $p_i$, both routes give $\sigma(p_i)$. The constant and connective
clauses agree, so structural induction proves commutativity.
Renaming changes the labels on the occurrences, while the pairs
$(r,t)$ continue to distinguish their positions. This permits the
same actual variable to occur in several inserted patterns without
identifying their separate occurrences.

For the substitution $p\mapsto A$ in
Proposition~\ref{prop:all-contexts}, the outer pattern is $p$,
so there is a single source block. This block is $\Occ(A)$.
The source and target formulas of the contextual equation are
$p\wedge\bot$ and $p\vee\top$; each has one occurrence of $p$. Their substituted occurrence sets therefore each
contain one copy of $\Occ(A)$, and the substituted arrow has type
$A\wedge\bot\to A\vee\top$. If the isolated literal is $\neg p$,
the substitution $p\mapsto\neg A$ instead gives
$\neg\neg A\wedge\bot\to\neg\neg A\vee\top$; the
double-negation square in Proposition~\ref{prop:all-contexts}
then gives the required equation for $A$.

\section{Diagrammatic forms of the structural equations}
\label{app:structural-diagrams}

\paragraph{The first-injection equation.}
For formulas $A,B,C$, equation~\eqref{eq:first-injection} is the
commutative triangle
\[
\begin{tikzcd}[column sep=large,row sep=large]
 A\wedge B
   \arrow[r,"1_A\wedge\iota_1^{B,C}"]
   \arrow[dr,"{\iota_1^{A\wedge B,C}}"']
   &A\wedge(B\vee C)\arrow[d,"d_{A,B,C}"]\\
   &(A\wedge B)\vee C.
\end{tikzcd}
\]
Both paths have source $A\wedge B$ and target $(A\wedge B)\vee C$.
The upper arrow is induced by $1_A:A\to A$ and
$\iota_1^{B,C}:B\to B\vee C$.

\paragraph{The zero identity.}
For a formula $A$, equation~\eqref{eq:zero-identity} is expressed by
the commutative diagram
\[
\begin{tikzcd}[column sep=large,row sep=large]
 A\arrow[r,"{\langle1_A,!_A\rangle}"]\arrow[d,"0_A"']
   &A\wedge\top\arrow[r,"1_A\wedge\eta_A"]
   &A\wedge(\neg A\vee A)\arrow[d,"d_{A,\neg A,A}"]\\
 A
   &\bot\vee A\arrow[l,"{[\kappa_A,1_A]}"']
   &(A\wedge\neg A)\vee A\arrow[l,"\epsilon_A\vee1_A"'].
\end{tikzcd}
\]
The path across the top, down the right side, and back along the
bottom is the five-arrow composite in~\eqref{eq:zero-identity}.
Its first arrow pairs $1_A:A\to A$ with $!_A:A\to\top$;
its last arrow copairs $\kappa_A:\bot\to A$ with $1_A:A\to A$.
The arrows between them use
$\eta_A:\top\to\neg A\vee A$ and
$\epsilon_A:A\wedge\neg A\to\bot$, together with the distributor.

\paragraph{The zero law.}
For any parallel arrows $f,g:A\to C$,
equation~\eqref{eq:zero-law} is the commutative square
\[
\begin{tikzcd}[column sep=large,row sep=large]
 A\arrow[r,"0_A"]\arrow[d,"g"']&A\arrow[d,"f"]\\
 C\arrow[r,"0_C"']&C.
\end{tikzcd}
\]
The upper-right path is $f\circ0_A:A\to C$, and the lower-left
path is $0_C\circ g:A\to C$. The arrows $f$ and $g$ need not be equal.

\paragraph{Use of AI tools.}
OpenAI's ChatGPT was used to assist with the development of the proof,
preparation of the manuscript, and checks of the arguments and references.

\begin{thebibliography}{9}
\small
\setlength{\itemsep}{2pt}
\setlength{\parskip}{0pt}
\bibitem{DP}
Kosta Do\v{s}en and Zoran Petri\'c,
\emph{Proof-Theoretical Coherence},
Studies in Logic, vol.~1, KCL Publications, London, 2004.
Revised author version, September 2007.
\url{https://www.mi.sanu.ac.rs/novi_sajt/biography/kosta/coh.pdf}.

\bibitem{DPCartesian}
Kosta Do\v{s}en and Zoran Petri\'c,
The maximality of cartesian categories,
\emph{Mathematical Logic Quarterly} \textbf{47} (2001), 137--144.
\url{https://arxiv.org/abs/math/9911059}.

\bibitem{DPBicartesian}
Kosta Do\v{s}en and Zoran Petri\'c,
Bicartesian coherence revisited,
in Zoran Ognjanovi\'c (ed.), \emph{Logic in Computer Science},
Zbornik radova \textbf{12(20)} (2009), 5--34.
\url{https://arxiv.org/abs/0711.4961}.
Author version: arXiv:0711.4961v3, December 8, 2008.

\bibitem{DPIsomorphic}
Kosta Do\v{s}en and Zoran Petri\'c,
Isomorphic formulae in classical propositional logic,
\emph{Mathematical Logic Quarterly} \textbf{58} (2012), 5--17.
\url{https://doi.org/10.1002/malq.201020020}.
Author version: \url{https://arxiv.org/abs/0912.1271},
arXiv:0912.1271v3, October 4, 2010.

\bibitem{DPSplit}
Kosta Do\v{s}en and Zoran Petri\'c,
Syntax for split preorders,
\emph{Annals of Pure and Applied Logic} \textbf{164} (2013), 443--481.
\url{https://arxiv.org/abs/0902.0742}.
Author version: arXiv:0902.0742v7, October 18, 2012.

\bibitem{LSFree}
Fran\c{c}ois Lamarche and Lutz Stra{\ss}burger,
Constructing free Boolean categories,
in \emph{Proceedings of the 20th Annual IEEE Symposium on Logic in
Computer Science (LICS 2005)}, IEEE Computer Society, 2005, pp.~209--218.
\url{https://doi.org/10.1109/LICS.2005.13}.

\bibitem{FP}
Carsten F\"uhrmann and David Pym,
Order-enriched categorical models of the classical sequent calculus,
\emph{Journal of Pure and Applied Algebra} \textbf{204} (2006), 21--78.
\url{https://doi.org/10.1016/j.jpaa.2005.03.016}.

\bibitem{Strassburger}
Lutz Stra{\ss}burger,
On the axiomatisation of Boolean categories with and without medial,
\emph{Theory and Applications of Categories} \textbf{18} (2007),
536--601.
\url{https://www.tac.mta.ca/tac/volumes/18/18/18-18abs.html}.

\bibitem{Notation}
Florian Lengyel,
\emph{The Unbearable Notation of Logic},
version~8, September 19, 2026, Section~9.
\url{https://github.com/flengyel/notes/blob/main/logic/The-Unbearable-Notation-of-Logic_v8.pdf}.
\end{thebibliography}
\end{document}
